\olchapter{pl}{com}{The Completeness Theorem}



\olfileid{pl}{com}{int}

\olsection{Introduction}

The completeness theorem is one of the most fundamental results about
logic.  It comes in two formulations, the equivalence of which we'll
prove.  In its first formulation it says something fundamental about
the relationship between semantic consequence and our !!{derivation}
system: if !!a{sentence}~$!A$ follows from some !!{sentence}s
$\Gamma$, then there is also !!a{derivation} that establishes $\Gamma
\Proves !A$. Thus, the !!{derivation} system is as strong as it can
possibly be without proving things that don't actually follow.

In its second formulation, it can be stated as a model existence
result: every consistent set of !!{sentence}s is satisfiable.
Consistency is a proof-theoretic notion: it says that our
!!{derivation} system is unable to produce certain !!{derivation}s.
But who's to say that just because there are no !!{derivation}s of a
certain sort from~$\Gamma$, it's guaranteed that there is
!!{valuation}{~$\pAssign{v}$}
with $\pSat{v}{\Gamma}$? Before the
completeness theorem was first proved---in fact before we had the
!!{derivation} systems we now do---the great German mathematician
David Hilbert held the view that consistency of mathematical theories
guarantees the existence of the objects they are about. He put it as
follows in a letter to Gottlob Frege:
\begin{quote}
  If the arbitrarily given axioms do not contradict one another with
  all their consequences, then they are true and the things defined by
  the axioms exist. This is for me the criterion of truth and
  existence. 
\end{quote}
Frege vehemently disagreed. The second formulation of the completeness
theorem shows that Hilbert was right in at least the sense that if the
axioms are consistent, then \emph{some}
!!{valuation} exists that makes them all
true.

These aren't the only reasons the completeness theorem---or rather,
its proof---is important.  It has a number of important consequences,
some of which we'll discuss separately.  For instance, since any
!!{derivation} that shows $\Gamma \Proves !A$ is finite and so can
only use finitely many of the !!{sentence}s in~$\Gamma$, it follows by
the completeness theorem that if $!A$ is a consequence of~$\Gamma$, it
is already a consequence of a finite subset of~$\Gamma$.  This is
called \emph{compactness}.  Equivalently, if every finite subset of
$\Gamma$ is consistent, then $\Gamma$ itself must be consistent.

Although the compactness theorem follows from the completeness theorem
via the detour through !!{derivation}s, it is also possible to use the
\emph{the proof of} the completeness theorem to establish it
directly. For what the proof does is take a set of !!{sentence}s with
a certain property---consistency---and constructs !!a{structure} out
of this set that has certain properties (in this case, that it
satisfies the set). Almost the very same construction can be used to
directly establish compactness, by starting from ``finitely
satisfiable'' sets of !!{sentence}s instead of consistent ones.






\olfileid{pl}{com}{out}

\olsection{Outline of the Proof}

The proof of the completeness theorem is a bit complex, and upon first
reading it, it is easy to get lost.  So let us outline the proof.  The
first step is a shift of perspective, that allows us to see a route to
a proof.  When completeness is thought of as ``whenever $\Gamma
\Entails !A$ then $\Gamma \Proves !A$,'' it may be hard to even come
up with an idea: for to show that $\Gamma \Proves !A$ we have to find
!!a{derivation}, and it does not look like the hypothesis that
$\Gamma \Entails !A$ helps us for this in any way.  For some proof
systems it is possible to directly construct !!a{derivation}, but we
will take a slightly different approach.  The shift in perspective
required is this: completeness can also be formulated as: ``if
$\Gamma$ is consistent, it is satisfiable.''  Perhaps we can use the
information in~$\Gamma$ together with the hypothesis that it is
consistent to construct !!a{valuation}
that satisfies every !!{formula}
in~$\Gamma$.  After all, we know
what kind of !!{valuation} we are looking
for: one that is as $\Gamma$ describes it!{}

If $\Gamma$ contains only !!{propositional variable}s, it is easy to construct
a model for it. All we have to do is come up with

  !!a{valuation}~$\pAssign{v}$ such that $\pSat{v}{p}$ for all $p \in
  \Gamma$. Well, let $\pAssign{v}(p) = \True$ iff $p \in \Gamma$.

Now suppose $\Gamma$ contains some !!{formula}~$\lnot !B$, with $!B$
atomic.  We might worry that the construction of
$\pAssign{v}$ interferes with the possibility
of making $\lnot !B$ true.  But here's where the consistency
of~$\Gamma$ comes in: if $\lnot !B \in \Gamma$, then $!B \notin
\Gamma$, or else $\Gamma$ would be inconsistent.  And if $!B \notin
\Gamma$, then according to our construction
of~$\pAssign{v}$,
$\pSat/{v}{!B}$, so
$\pSat{v}{\lnot !B}$.  So far so good.

What if $\Gamma$ contains complex, non-atomic formulas? Say it
contains $!A \land !B$. To make that true, we should proceed as if
both $!A$ and $!B$ were in~$\Gamma$.  And if $!A \lor !B \in \Gamma$,
then we will have to make at least one of them true, i.e., proceed as
if one of them was in~$\Gamma$.

This suggests the following idea: we add additional !!{formula}s
to~$\Gamma$ so as to (a)~keep the resulting set consistent and
(b)~make sure that for every possible atomic !!{sentence}~$!A$, either
$!A$ is in the resulting set, or $\lnot !A$ is, and (c)~such that,
whenever $!A \land !B$ is in the set, so are both $!A$ and $!B$, if
$!A \lor !B$ is in the set, at least one of $!A$ or $!B$ is also, etc.
We keep doing this (potentially forever).  Call the set of all
!!{formula}s so added~$\Gamma^*$.  Then our construction above would
provide us with !!a{valuation}~$\pAssign{v}$
for which we could prove, by induction, that it satisfies all sentences
in~$\Gamma^*$, and hence also all sentence in~$\Gamma$
since $\Gamma \subseteq \Gamma^*$.  It turns out that guaranteeing (a)
and~(b) is enough. A set of sentences for which (b) holds is called
\emph{complete}. So our task will be to extend the consistent
set~$\Gamma$ to a consistent and complete set~$\Gamma^*$.



So here's what we'll do. First we investigate the properties of
!!{complete} consistent sets, in particular we prove that
!!a{complete} consistent set contains $!A \land !B$ iff it contains
both $!A$ and~$!B$, $!A \lor !B$ iff it contains at least one of them,
etc. (\olref[ccs]{prop:ccs}).
We'll then take the consistent
set~$\Gamma$ and show that it can be extended
to a consistent and
!!{complete} set~$\Gamma^*$ (\olref[lin]{lem:lindenbaum}). This set
$\Gamma^*$ is what we'll use to define our !!{valuation}~$\pAssign v(\Gamma^*)$.
The valuation is determined by the !!{propositional
    variable}s in~$\Gamma^*$ (\olref[mod]{defn:termmodel}).  We'll
use the properties of complete consistent
sets to show that indeed
$\pSat{v(\Gamma^*)}{!A}$ iff $!A \in
\Gamma^*$ (\olref[mod]{lem:truth}), and thus in particular,
$\pSat{v(\Gamma^*)}{\Gamma}$.






\olfileid{pl}{com}{ccs}
      
\olsection{Complete Consistent Sets of \usetoken{P}{sentence}}

\begin{defn}[Complete set]
\ollabel{def:complete-set} A set~$\Gamma$ of !!{sentence}s is
\emph{!!{complete}} iff for any !!{sentence}~$!A$, either $!A \in
\Gamma$ or $\lnot !A \in \Gamma$.
\end{defn}

\begin{explain}
!!^{complete} sets of sentences leave no questions unanswered. For
any !!{sentence}~$!A$, $\Gamma$ ``says'' if $!A$ is true or false.  The
importance of !!{complete} sets extends beyond the proof of the
completeness theorem. A theory which is !!{complete} and
axiomatizable, for instance, is always decidable.
\end{explain}

\begin{explain}
!!^{complete} consistent sets are important in the completeness proof
since we can guarantee that every consistent set of
!!{sentence}s~$\Gamma$ is contained in !!a{complete} consistent
set~$\Gamma^*$.  !!^a{complete} consistent set contains, for each
!!{sentence}~$!A$, either $!A$ or its negation $\lnot !A$, but not
both. This is true in particular for propositional variables, so from !!a{complete}
consistent set, we can construct
!!a{valuation} where the
truth value assigned to
propositional variables is defined according to which
!!{propositional variable}s are
in~$\Gamma^*$. This !!{valuation} can then
be shown to make all !!{sentence}s in~$\Gamma^*$ (and hence also all
those in~$\Gamma$) true. The proof of this latter fact requires that
$\lnot !A \in \Gamma^*$ iff $!A \notin \Gamma^*$, $(!A \lor !B) \in
\Gamma^*$ iff $!A \in \Gamma^*$ or $!B \in \Gamma^*$, etc.
\end{explain}

In what follows, we will often tacitly use the properties of
reflexivity, monotonicity, and transitivity of $\Proves$ (see

  \tagrefs{prfSC/{pl:seq:ptn:sec},prfND/{pl:ntd:ptn:sec},prfAX/{pl:axd:ptn:sec},prfTab/{pl:tab:ptn:sec}}).

\begin{prop}
\ollabel{prop:ccs}
Suppose $\Gamma$ is !!{complete} and consistent. Then:
\begin{enumerate}
\item \ollabel{prop:ccs-prov-in} If $\Gamma \Proves !A$, then $!A \in
  \Gamma$.

\ollabel{prop:ccs-and} $!A \land !B \in \Gamma$
  iff both $!A \in \Gamma$ and $!B \in \Gamma$.

\ollabel{prop:ccs-or} $!A \lor !B \in \Gamma$ iff
  either $!A \in \Gamma$ or $!B \in \Gamma$.

\ollabel{prop:ccs-if} $!A \lif !B \in \Gamma$ iff
  either $!A \notin \Gamma$ or $!B \in \Gamma$.
\end{enumerate}
\end{prop}

\begin{proof}
Let us suppose for all of the following that $\Gamma$ is !!{complete} and
consistent.
\begin{enumerate}
\item If $\Gamma \Proves !A$, then $!A \in \Gamma$.

Suppose that $\Gamma \Proves !A$. Suppose to the contrary that $!A
\notin \Gamma$. Since $\Gamma$ is !!{complete}, $\lnot !A \in \Gamma$.
By

  \tagrefs{prfSC/{pl:seq:prv:prop:explicit-inc},
    prfND/{pl:ntd:prv:prop:explicit-inc},
    prfAX/{pl:axd:prv:prop:explicit-inc},
    prfTab/{pl:tab:prv:prop:explicit-inc}},
$\Gamma$ is inconsistent. This contradicts the assumption that
$\Gamma$ is consistent. Hence, it cannot be the case that $!A \notin
\Gamma$, so $!A \in \Gamma$.


Exercise.



First we show that if $!A \lor !B \in \Gamma$, then either $!A \in
\Gamma$ or $!B \in \Gamma$. Suppose $!A \lor !B \in \Gamma$ but $!A
\notin \Gamma$ and $!B \notin \Gamma$.  Since $\Gamma$ is
!!{complete}, $\lnot !A \in \Gamma$ and $\lnot !B \in \Gamma$. By

  \tagrefs{prfSC/{pl:seq:ppr:prop:provability-lor},
    prfND/{pl:ntd:ppr:prop:provability-lor},
    prfAX/{pl:axd:ppr:prop:provability-lor},
    prfTab/{pl:tab:ppr:prop:provability-lor}},
item (1), $\Gamma$ is inconsistent, a contradiction. Hence, either $!A
\in \Gamma$ or $!B \in \Gamma$.

For the reverse direction, suppose that $!A \in \Gamma$ or $!B \in
\Gamma$. By

  \tagrefs{prfSC/{pl:seq:ppr:prop:provability-lor},
    prfND/{pl:ntd:ppr:prop:provability-lor},
    prfAX/{pl:axd:ppr:prop:provability-lor},
    prfTab/{pl:tab:ppr:prop:provability-lor}}, item (2),
$\Gamma \Proves !A \lor !B$. By \olref{prop:ccs-prov-in}, $!A \lor
!B \in \Gamma$, as required.


Exercise.
\end{enumerate}
\end{proof}




\begin{prob}
Complete the proof of \olref[pl][com][ccs]{prop:ccs}.
\end{prob}








\olfileid{pl}{com}{lin}

\olsection{Lindenbaum's Lemma}

\begin{explain}
We now prove a lemma that shows that any consistent set of
!!{sentence}s is contained in some set of sentences which is not just
consistent, but also !!{complete}. The proof works by adding one
!!{sentence} at a time, guaranteeing at each step that the set remains
consistent. We do this so that for every $!A$, either $!A$ or $\lnot
!A$ gets added at some stage. The union of all stages in that
construction then contains either $!A$ or its negation~$\lnot !A$ and
is thus complete. It is also consistent, since we make sure at each
stage not to introduce an inconsistency.
\end{explain}

\begin{lem}[Lindenbaum's Lemma]
\ollabel{lem:lindenbaum} Every consistent
set~$\Gamma$ in a language~$\Lang{L}$ can be
extended to !!a{complete} and consistent set~$\Gamma^*$.
\end{lem}

\begin{proof}
Let $\Gamma$ be consistent.  Let $!A_0$, $!A_1$,
\dots{} be an enumeration of all the !!{sentence}s of~$\Lang L$.
Define $\Gamma_0 = \Gamma$, and
\[
\Gamma_{n+1} =
\begin{cases}
\Gamma_n \cup \{ !A_n \} & \textrm{if $\Gamma_n \cup \{!A_n\}$ is
  consistent;} \\
\Gamma_n \cup \{ \lnot !A_n \} & \textrm{otherwise.}
\end{cases}
\]
Let $\Gamma^* = \bigcup_{n \geq 0} \Gamma_n$.

Each $\Gamma_n$ is consistent: $\Gamma_0$ is consistent by definition.
If $\Gamma_{n+1} = \Gamma_n \cup \{!A_n\}$, this is because the latter
is consistent.  If it isn't, $\Gamma_{n+1} = \Gamma_n \cup \{\lnot
!A_n\}$. We have to verify that $\Gamma_n \cup \{\lnot !A_n\}$ is
consistent. Suppose it's not. Then \emph{both} $\Gamma_n \cup
\{!A_n\}$ and $\Gamma_n \cup \{\lnot !A_n\}$ are inconsistent.  This
means that $\Gamma_n$ would be inconsistent by

  \tagrefs{prfAX/{pl:axd:prv:prop:provability-exhaustive},
    prfSC/{pl:seq:prv:prop:provability-exhaustive},
    prfND/{pl:ntd:prv:prop:provability-exhaustive},
    prfTab/{pl:tab:prv:prop:provability-exhaustive}},
contrary to the induction hypothesis.

For every~$n$ and every $i < n$, $\Gamma_i \subseteq \Gamma_n$. This
follows by a simple induction on~$n$. For $n=0$, there are no $i < 0$,
so the claim holds automatically.  For the inductive step, suppose it
is true for~$n$. We show that if $i < n+1$ then $\Gamma_i \subseteq
\Gamma_{n+1}$. We have $\Gamma_{n+1} = \Gamma_n \cup \{!A_n\}$ or $=
\Gamma_n \cup \{\lnot !A_n\}$ by construction. So $\Gamma_n \subseteq
\Gamma_{n+1}$. If $i < n+1$, then $\Gamma_i \subseteq \Gamma_n$ by
inductive hypothesis (if $i < n$) or the trivial fact that $\Gamma_n
\subseteq \Gamma_n$ (if $i = n$). We get that $\Gamma_i \subseteq
\Gamma_{n+1}$ by transitivity of~$\subseteq$.

From this it follows that $\Gamma^*$ is consistent. Here's why: Let
$\Gamma' \subseteq \Gamma^*$ be finite. Each $!B \in \Gamma'$ is also
in~$\Gamma_i$ for some~$i$. Let $n$ be the largest of these. Since
$\Gamma_i \subseteq \Gamma_n$ if $i \le n$, every $!B \in \Gamma'$ is
also $\in \Gamma_n$, i.e., $\Gamma' \subseteq \Gamma_n$, and
$\Gamma_n$~is consistent. So, every finite subset $\Gamma' \subseteq
\Gamma^*$ is consistent. By 
  \tagrefs{prfAX/{pl:axd:ptn:prop:proves-compact},
    prfSC/{pl:seq:ptn:prop:proves-compact},
    prfND/{pl:ntd:ptn:prop:proves-compact},
    prfTab/{pl:tab:ptn:prop:proves-compact}}, $\Gamma^*$ is
  consistent.

Every !!{sentence} of $\Frm[L]$ appears on the list used to
define~$\Gamma^*$. If $!A_n \notin \Gamma^*$, then that is because
$\Gamma_n \cup \{!A_n\}$ was inconsistent.  But then $\lnot !A_n
\in \Gamma^*$, so $\Gamma^*$ is !!{complete}.
\end{proof}





\olfileid{pl}{com}{mod}

\olsection{Construction of a Model}

\begin{explain}
We are now ready to define !!a{valuation}
  that makes all $!A \in \Gamma$ true. To do this, we first apply
  Lindenbaum's Lemma: we get a complete consistent $\Gamma^* \supseteq
  \Gamma$. We let the !!{propositional variable}s in~$\Gamma^*$
  determine~$\pAssign v(\Gamma^*)$.
\end{explain}


\begin{defn}
\ollabel{defn:termmodel}
Suppose $\Gamma^*$ is a complete consistent set of !!{formula}s. Then we let
\[
\pAssign v(\Gamma^*)(p) = \begin{cases}
  \True & \text{if $p \in \Gamma^*$}\\
  \False & \text{if $p \notin \Gamma^*$}
  \end{cases}
\]
\end{defn}


\iftag{FOL}{
\begin{explain}
  !!^a{structure}~$\Struct{M}$ may make an existentially
    quantified !!{sentence}~$\lexists[x][!A(x)]$ true without there
    being an instance~$!A(t)$ that it makes true.
  !!^a{structure}~$\Struct{M}$ may make all
    instances~$!A(t)$ of a universally quantified
    !!{sentence}~$\lforall[x][!A(x)]$ true, without
    making~$\lforall[x][!A(x)]$ true. This is because in general
  not every !!{element} of~$\Domain{M}$ is the value of a closed term
  ($\Struct{M}$ may not be covered).  This is the reason the
  satisfaction relation is defined via variable assignments. However,
  for our term model~$\Struct{M(\Gamma^*)}$ this wouldn't be
  necessary---because it is covered. This is the content of the next
  result.
\end{explain}

\begin{prop}
\ollabel{prop:quant-termmodel}
Let $\Struct M(\Gamma^*)$ be the term model of \olref{defn:termmodel}.
\begin{tagenumerate}{prvEx,prvAll}
$\Sat{M(\Gamma^*)}{\lexists[x][!A(x)]}$ iff
  $\Sat{M(\Gamma^*)}{!A(t)}$ for at least one closed term~$t$.
$\Sat{M(\Gamma^*)}{\lforall[x][!A(x)]}$ iff
  $\Sat{M(\Gamma^*)}{!A(t)}$ for all closed terms~$t$.
\end{tagenumerate}
\end{prop}

\begin{proof}
\begin{tagenumerate}{prvEx,prvAll}

  By \olref[syn][ass]{prop:sat-quant},
    $\Sat{M(\Gamma^*)}{\lexists[x][!A(x)]}$ iff for at least one
    variable assignment~$s$, $\Sat{M(\Gamma^*)}{!A(x)}[s]$. As
    $\Domain{M(\Gamma^*)}$ consists of the closed terms of~$\Lang{L}$,
    this is the case iff there is at least one closed term~$t$ such
    that $s(x) = t$ and $\Sat{M(\Gamma^*)}{!A(x)}[s]$.  By
    \olref[fol][syn][ext]{prop:ext-formulas},
    $\Sat{M(\Gamma^*)}{!A(x)}[s]$ iff $\Sat{M(\Gamma^*)}{!A(t)}[s]$,
    where $s(x) = t$.  By \olref[fol][syn][ass]{prop:sentence-sat-true},
    $\Sat{M(\Gamma^*)}{!A(t)}[s]$ iff $\Sat{M(\Gamma^*)}{!A(t)}$,
    since $!A(t)$ is a sentence.

  Exercise.
\end{tagenumerate}
\end{proof}
}{}



\begin{lem}[Truth Lemma]
\ollabel{lem:truth} 
$\pSat{v(\Gamma^*)}{!A}$ iff $!A \in \Gamma^*$.
\end{lem}

\begin{proof}
We prove both directions simultaneously, and by induction on $!A$.
\begin{enumerate}
\indcase{!A}{\lfalse}
  {$\pSat/{v(\Gamma^*)}{\lfalse}$
    by definition of satisfaction. On the other hand, $\lfalse \notin
    \Gamma^*$ since $\Gamma^*$ is consistent.}



\indcase{!A}{p}{$\pSat{v(\Gamma^*)}{p}$ iff
    $\pAssign v(\Gamma^*)(p) = \True$ (by the definition of
    satisfaction) iff $p \in \Gamma^*$ (by the construction of
    $\pAssign v(\Gamma^*)$).}


  \indcase{!A}{\lnot !B}
    {$\pSat{v(\Gamma^*)}{\indfrm}$ iff
    $\pSat/{v(\Gamma^*)}{!B}$ (by
    definition of satisfaction). By induction hypothesis,
    $\pSat/{v(\Gamma^*)}{!B}$ iff
    $!B \notin \Gamma^*$. Since $\Gamma^*$ is consistent and
    !!{complete}, $!B
    \notin \Gamma^*$ iff $\lnot !B \in \Gamma^*$.}



  \indcase!{!A}{!B \land !C}{}



  \indcase{!A}{!B \lor !C}
    {$\pSat{v(\Gamma^*)}{\indfrm}$
    iff $\pSat{v(\Gamma^*)}{!B}$ or
    $\pSat{v(\Gamma^*)}{!C}$ (by
    definition of satisfaction) iff $!B \in \Gamma^*$ or $!C \in
    \Gamma^*$ (by induction hypothesis). This is the case iff $(!B
    \lor !C) \in \Gamma^*$ (by
    \olref[ccs]{prop:ccs}\olref[ccs]{prop:ccs-or}).}
  


  \indcase!{!A}{!B \lif !C}{}


\end{enumerate}
\end{proof}






\begin{prob}
Complete the proof of \olref[pl][com][mod]{lem:truth}.
\end{prob}









\olfileid{pl}{com}{cth}

\olsection{The Completeness Theorem}

\begin{explain}
Let's combine our results: we arrive at the completeness theorem.
\end{explain}

\begin{thm}[Completeness Theorem]
\ollabel{thm:completeness}
Let $\Gamma$ be a set of !!{sentence}s.  If $\Gamma$ is consistent, it
is satisfiable.
\end{thm}

\begin{proof}
Suppose $\Gamma$ is consistent. By \olref[lin]{lem:lindenbaum}, there
is a $\Gamma^* \supseteq \Gamma$ which is
consistent and !!{complete}.
By
\olref[mod]{lem:truth},
$\pSat{v(\Gamma^*)}{!A}$ iff $!A \in
\Gamma^*$.  From this it follows in particular that for all $!A \in
\Gamma$, $\pSat{v(\Gamma^*)}{!A}$, so
$\Gamma$ is satisfiable.
\end{proof}

\begin{cor}[Completeness Theorem, Second Version]
\ollabel{cor:completeness}
For all $\Gamma$ and !!{sentence}s~$!A$: if $\Gamma \Entails !A$ then
$\Gamma \Proves !A$.
\end{cor}

\begin{proof}
Note that the $\Gamma$'s in \olref{cor:completeness} and
\olref{thm:completeness} are universally quantified.  To make sure we
do not confuse ourselves, let us restate \olref{thm:completeness}
using a different variable: for any set of !!{sentence}s~$\Delta$, if
$\Delta$ is consistent, it is satisfiable.  By contraposition, if
$\Delta$ is not satisfiable, then $\Delta$ is inconsistent.  We will
use this to prove the corollary.

Suppose that $\Gamma \Entails !A$.  Then $\Gamma \cup \{\lnot !A\}$ is
unsatisfiable by \olref[syn][sem]{prop:entails-unsat}.  Taking $\Gamma
\cup \{\lnot !A\}$ as our $\Delta$, the previous version of
\olref{thm:completeness} gives us that $\Gamma \cup \{\lnot !A\}$ is
inconsistent.  By

  \tagrefs{prfAX/{pl:axd:prv:prop:prov-incons},
    prfSC/{pl:seq:prv:prop:prov-incons},
    prfND/{pl:ntd:prv:prop:prov-incons},
    prfTab/{pl:tab:prv:prop:prov-incons}},
$\Gamma \Proves !A$.
\end{proof}




\begin{prob}
Use \olref[pl][com][cth]{cor:completeness} to prove
\olref[pl][com][cth]{thm:completeness}, thus showing that the two
formulations of the completeness theorem are equivalent.
\end{prob}





\begin{prob}
In order for !!a{derivation} system to be complete, its rules must be
strong enough to prove every unsatisfiable set inconsistent.  Which of
the rules of !!{derivation} were necessary to prove completeness?  Are any
of these rules not used anywhere in the proof?  In order to answer
these questions, make a list or diagram that shows which of the rules
of !!{derivation} were used in which results that lead up to the proof of
\olref[pl][com][cth]{thm:completeness}.  Be sure to note any tacit
uses of rules in these proofs.
\end{prob}






\olfileid{pl}{com}{com}

\olsection{The Compactness Theorem}

One important consequence of the completeness theorem is the
compactness theorem.  The compactness theorem states that if each
\emph{finite} subset of a set of !!{sentence}s is satisfiable, the
entire set is satisfiable---even if the set itself is infinite. This
is far from obvious. There is nothing that seems to rule out, at first
glance at least, the possibility of there being infinite sets of
!!{sentence}s which are contradictory, but the contradiction only
arises, so to speak, from the infinite number.  The compactness
theorem says that such a scenario can be ruled out: there are no
unsatisfiable infinite sets of !!{sentence}s each finite subset of
which is satisfiable. Like the completeness theorem, it has a version
related to entailment: if an infinite set of !!{sentence}s entails
something, already a finite subset does.

\begin{defn}
  A set $\Gamma$ of !!{formula}s is \emph{finitely satisfiable} iff every finite $\Gamma_0 \subseteq \Gamma$ is satisfiable.
\end{defn}

\begin{thm}[Compactness Theorem]
\ollabel{thm:compactness}
The following hold for any sentences $\Gamma$ and $!A$:
\begin{enumerate}
  \item $\Gamma \Entails !A$ iff there is a finite $\Gamma_0
    \subseteq \Gamma$ such that $\Gamma_0 \Entails !A$.
  \item $\Gamma$ is satisfiable iff it is finitely
    satisfiable.
\end{enumerate}
\end{thm}

\begin{proof}
We prove (2).  If $\Gamma$ is satisfiable, then there is
!!a{valuation}~$\pAssign{v}$
such that $\pSat{v}{!A}$ for all $!A \in
\Gamma$.  Of course, this $\pAssign{v}$ also
satisfies every finite subset of~$\Gamma$, so $\Gamma$ is finitely
satisfiable.

Now suppose that $\Gamma$ is finitely satisfiable.  Then every finite
subset~$\Gamma_0 \subseteq \Gamma$ is satisfiable.  By soundness
(
  \tagrefs{prfAX/{pl:axd:sou:cor:consistency-soundness},
    prfSC/{pl:seq:sou:cor:consistency-soundness},
    prfND/{pl:ntd:sou:cor:consistency-soundness},
    prfTab/{pl:tab:sou:cor:consistency-soundness}}),
every finite subset is consistent.  Then $\Gamma$ itself must be
consistent by

  \tagrefs{prfAX/{pl:axd:ptn:prop:proves-compact},
    prfSC/{pl:seq:ptn:prop:proves-compact},
    prfND/{pl:ntd:ptn:prop:proves-compact},
    prfTab/{pl:tab:ptn:prop:proves-compact}}.
By completeness (\olref[cth]{thm:completeness}), since $\Gamma$~is
consistent, it is satisfiable.
\end{proof}




\begin{prob}
Prove (1) of \olref[pl][com][com]{thm:compactness}.
\end{prob}












\olfileid{pl}{com}{cpd}

\olsection{A Direct Proof of the Compactness Theorem}

We can prove the Compactness Theorem directly, without appealing to
the Completeness Theorem, using the same ideas as in the proof of the
completeness theorem.  In the proof of the Completeness Theorem we
started with a consistent set~$\Gamma$ of !!{sentence}s, expanded it
to a consistent and !!{complete}
set~$\Gamma^*$ of !!{sentence}s, and then showed that in the
!!{valuation}~$\pAssign{v(\Gamma^*)}$
constructed from $\Gamma^*$, all !!{sentence}s of~$\Gamma$ are true,
so $\Gamma$ is satisfiable.

We can use the same method to show that a finitely satisfiable set of
sentences is satisfiable. We just have to prove the corresponding
versions of the results leading to the truth lemma where we replace
``consistent'' with ``finitely satisfiable.''

\begin{prop}
\ollabel{prop:fsat-ccs}
Suppose $\Gamma$ is !!{complete} and finitely satisfiable. Then:
\begin{enumerate}
$(!A \land !B) \in \Gamma$
  iff both $!A \in \Gamma$ and $!B \in \Gamma$.

$(!A \lor !B) \in \Gamma$ iff
  either $!A \in \Gamma$ or $!B \in \Gamma$.

$(!A \lif !B) \in \Gamma$ iff
  either $!A \notin \Gamma$ or $!B \in \Gamma$.
\end{enumerate}
\end{prop}




\begin{prob}
Prove \olref[pl][com][cpd]{prop:fsat-ccs}. Avoid the use of $\Proves$.
\end{prob}










\begin{lem}
\ollabel{lem:fsat-lindenbaum} Every finitely satisfiable set~$\Gamma$
can be extended to !!a{complete} and finitely satisfiable
set~$\Gamma^*$.
\end{lem}




\begin{prob}
Prove \olref[pl][com][cpd]{lem:fsat-lindenbaum}.  (Hint: the crucial
step is to show that if $\Gamma_n$ is finitely satisfiable, then
either $\Gamma_n \cup \{!A_n\}$ or $\Gamma_n \cup \{\lnot !A_n\}$ is
finitely satisfiable.)
\end{prob}


\begin{thm}[Compactness]
\ollabel{thm:compactness-direct} $\Gamma$ is satisfiable if and only
if it is finitely satisfiable.
\end{thm}

\begin{proof}
If $\Gamma$ is satisfiable, then there is
!!a{valuation}~$\pAssign{v}$
such that $\pSat{v}{!A}$ for all $!A \in \Gamma$.
Of course, this $\pAssign{v}$ also
satisfies every finite subset of~$\Gamma$, so $\Gamma$ is finitely
satisfiable.

Now suppose that $\Gamma$ is finitely satisfiable. By \olref{lem:fsat-lindenbaum},
$\Gamma$ can be extended to !!a{complete} and
finitely satisfiable set~$\Gamma^*$. Construct the !!{valuation}~$\pAssign{v(\Gamma^*)}$
as in \olref[mod]{defn:termmodel}.
The proof of the Truth Lemma (\olref[mod]{lem:truth}) goes through if
we replace references to \olref[ccs]{prop:ccs}.
\end{proof}




\begin{prob}
Write out the complete proof of the Truth Lemma
(\olref[pl][com][mod]{lem:truth}) in the version required for the
proof of \olref[pl][com][cpd]{thm:compactness-direct}.
\end{prob}






\OLEndChapterHook





\OLEndPartHook

